\documentclass[12pt,a4paper,oneside]{report}
\usepackage[english]{../../../../latex-setup/learningnotes}
% The shared style is made for the notes; a cheatsheet needs narrower
% margins and no running headers.
\geometry{margin=18mm}
\pagestyle{notescontents}
% Short boxes read better when a page break never splits them.
\tcbset{notesbox/.append style={unbreakable}}
% Each heading starts the next chapter of the notes, so the box
% counters restart and every number matches the notes.
\newcommand{\notesheading}[1]{%
  \refstepcounter{chapter}%
  \section*{\chaptername~\thechapter.\enspace #1}}
\begin{document}
\thispagestyle{notescontents}
\begin{center}
  {\Large\bfseries Fundamentals of Statistical Learning}\par
  \smallskip
  {\large Definitions and results cheatsheet}
\end{center}

\notesheading{Defining probability}

\begin{frameddefn}{Outcomes and events}\label{def:outcomes-events}
A random experiment is a process whose result is not known in advance.
Each possible result is an \emph{outcome} $\omega$. The set of all possible
outcomes is the \emph{sample space} $\Omega$. An \emph{event} is a subset
$A\subseteq\Omega$; it occurs when the observed outcome belongs to $A$.
\end{frameddefn}

\begin{framedprop}{De Morgan's laws}\label{prop:de-morgan}
For any events $A$ and $B$,
\[
  (A\cup B)^c=A^c\cap B^c,
  \qquad
  (A\cap B)^c=A^c\cup B^c.
\]
More generally, for events $A_1,\ldots,A_n$,
\[
  \left(\bigcup_{i=1}^{n}A_i\right)^c
    =\bigcap_{i=1}^{n}A_i^c,
  \qquad
  \left(\bigcap_{i=1}^{n}A_i\right)^c
    =\bigcup_{i=1}^{n}A_i^c.
\]
\end{framedprop}

\begin{frameddefn}{Probability on a finite sample space}\label{def:probability-finite}
A probability measure $P$ assigns a number $P(A)$ to every event
$A\subseteq\Omega$ such that:
\begin{enumerate}
  \item $0\leq P(A)\leq 1$ for every event $A$;
  \item $P(\Omega)=1$;
  \item if $A\cap B=\varnothing$, then
        $P(A\cup B)=P(A)+P(B)$.
\end{enumerate}
The last property is called \emph{finite additivity}.
\end{frameddefn}

\begin{framedprop}{Inclusion--exclusion for two events}
For any events $A$ and $B$,
\[
  P(A\cup B)=P(A)+P(B)-P(A\cap B).
\]
\end{framedprop}

\begin{framedprop}{Union bound (Boole's inequality)}
For events $A_1,\ldots,A_n$,
\[
  P\!\left(\bigcup_{i=1}^{n} A_i\right)
  \leq \sum_{i=1}^{n}P(A_i).
\]
\end{framedprop}

\begin{framedprop}{Bonferroni's inequality}
For events $A_1,\ldots,A_n$,
\[
  P\!\left(\bigcap_{i=1}^{n} A_i\right)
  \geq \sum_{i=1}^{n}P(A_i)-(n-1).
\]
\end{framedprop}

\notesheading{Conditional probability}

\begin{frameddefn}{Conditional probability}
\label{def:conditional-probability}
For events $A,B$ with $P(B)>0$, the probability of $A$ given $B$ is
\[
  P(A\mid B)=\frac{P(A\cap B)}{P(B)}.
\]
\end{frameddefn}

\begin{frameddefn}{Independent events}
\label{def:independent-events}
Two events $A$ and $B$ are independent if
\[
  P(A\cap B)=P(A)P(B).
\]
\end{frameddefn}

\begingroup
\renewcommand{\thefootnote}{\fnsymbol{footnote}}
\footnotetext[2]{Here $B_1,\ldots,B_m$ form a partition of the sample space, with $P(B_i)>0$.}
\endgroup
\nopagebreak
\begin{framedthm}{Bayes' theorem}
\label{thm:bayes}
For a group $B_j$ of the partition\begingroup\renewcommand{\thefootnote}{\fnsymbol{footnote}}\footnotemark[2]\endgroup\ and an event $A$ with $P(A)>0$,
\[
  P(B_j\mid A)
    =\frac{P(A\mid B_j)P(B_j)}{P(A)}
    =\frac{P(A\mid B_j)P(B_j)}
           {\sum_{i=1}^m P(A\mid B_i)P(B_i)}.
\]
\end{framedthm}

\begin{frameddefn}{Prior and posterior}
In Bayes' theorem, $P(B_j)$ is the probability of $B_j$ before
observing $A$, called the \textbf{prior}. Its update $P(B_j\mid A)$
after observing $A$ is called the \textbf{posterior}.
\end{frameddefn}

\notesheading{Random variables}

\begingroup
\renewcommand{\thefootnote}{\fnsymbol{footnote}}
\footnotetext[1]{A rigorous definition of measurability is beyond the scope of this course.}
\endgroup
\nopagebreak
\begin{frameddefn}{Random object}\label{def:random-object}
A \emph{random object} is a measurable\begingroup\renewcommand{\thefootnote}{\fnsymbol{footnote}}\footnotemark[1]\endgroup\ map from outcomes to possible
observations:
\[
  Z:\Omega\longrightarrow\mathcal T,
  \qquad \omega\longmapsto Z(\omega).
\]
The set $\mathcal T$ says what kind of thing we record. It could contain
numbers, vectors, images, or text messages. Once $\omega$ is known,
$Z(\omega)$ is the observation we get.
\end{frameddefn}

\begin{frameddefn}{Random variable}\label{def:random-variable}
A \emph{random variable} is a measurable map
\[
  X:\Omega\longrightarrow\mathbb R.
\]
It assigns one real number $X(\omega)$ to each outcome $\omega$.
\end{frameddefn}

\begin{frameddefn}{Random vector}\label{def:random-vector}
A \emph{$p$-dimensional random vector} is a measurable map
\[
  \underline{X}:\Omega\longrightarrow\mathbb R^p,
  \qquad
  \underline{X}(\omega)
    :=\bigl(X_1(\omega),\ldots,X_p(\omega)\bigr).
\]
Each coordinate $X_j$ is a real-valued random variable. Thus one
outcome gives us $p$ numbers collected into a vector. These measurements
can be related; they do not have to be independent.
\end{frameddefn}

\begin{frameddefn}{Preimage}\label{def:preimage}
For a random variable $X:\Omega\to\mathbb R$ and a set
$E\subseteq\mathbb R$, the \emph{preimage} (or \emph{inverse image}) of
$E$ is
\[
  X^{-1}(E):=\{\omega\in\Omega:X(\omega)\in E\}\subseteq\Omega.
\]
To find this event, we collect the outcomes whose outputs fall in $E$.
Whenever $X^{-1}(E)$ is an event, we define
\[
  \begin{aligned}
    P(X\in E)&:=P\bigl(X^{-1}(E)\bigr)\\
              &=P\bigl(\{\omega\in\Omega:X(\omega)\in E\}\bigr).
  \end{aligned}
\]
In particular, for a single value $x$,
\[
  P(X=x):=P\bigl(X^{-1}(\{x\})\bigr).
\]
Here $X$ is the random variable, while $x$ is a particular real number.
\end{frameddefn}

\begin{frameddefn}{Distribution of a random variable}
\label{def:distribution}
The \emph{distribution} of $X$ is the rule that tells us with what
probability $X$ takes values in a given set $E$.
\[
  P_X(E)=P(X\in E).
\]
\end{frameddefn}

\begingroup
\renewcommand{\thefootnote}{\fnsymbol{footnote}}
\footnotetext[3]{A finite set has finitely many elements. A set $A$ is countably infinite if there exists a bijection $f:\mathbb N\to A$.}
\endgroup
\nopagebreak
\begin{frameddefn}{Discrete random variable and PMF}
\label{def:discrete-pmf}
A random variable $X$ is \emph{discrete} if there is a finite or
countably infinite set\begingroup\renewcommand{\thefootnote}{\fnsymbol{footnote}}\footnotemark[3]\endgroup\ 
$\mathcal X=\{x_1,x_2,\ldots\}$ such that
$P(X\in\mathcal X)=1$.

Its \emph{probability mass function} (PMF) is
\[
  p_X(x):=P(X=x),\qquad x\in\mathbb R.
\]
\begin{enumerate}
  \item $p_X(x)\geq0$ for every $x\in\mathbb R$.
  \item The \emph{support}, in the discrete sense used here, is
    $\operatorname{supp}(X):=\{x\in\mathbb R:p_X(x)>0\}$.
    Thus the support lists the values that have positive probability;
    values outside it have probability zero.
  \item $\displaystyle\sum_{x\in\mathcal X}p_X(x)=1$.
\end{enumerate}
\end{frameddefn}

\begin{frameddefn}{Continuous random variable and PDF}
\label{def:continuous-pdf}
A random variable $X$ is \emph{continuous} (in the sense used here)
if there exists a function $f_X$, called its \emph{probability density
function} (PDF), such that, for every $a\leq b$,
\[
  P(a\leq X\leq b)=\int_a^b f_X(x)\,dx.
\]
\begin{enumerate}
  \item $f_X(x)\geq0$ for every $x\in\mathbb R$.
  \item Following the convention in the slides, the \emph{support} is
    $\operatorname{supp}(X):=\{x\in\mathbb R:f_X(x)>0\}$.
  \item $\displaystyle\int_{\mathbb R}f_X(x)\,dx=1$, since the total
    area under $f_X$ equals $P(X\in\mathbb R)$, and $X$ certainly
    takes some real value.
\end{enumerate}
\end{frameddefn}

\begin{frameddefn}{Cumulative distribution function}
\label{def:cdf}
Given a random variable $X$, its \emph{cumulative distribution
function} (CDF) is the function $F_X:\mathbb R\to[0,1]$ defined by
\[
  F_X(x):=P(X\leq x)=P\bigl(X\in(-\infty,x]\bigr),
  \qquad x\in\mathbb R.
\]
It gives the probability accumulated up to and including the
threshold $x$. It is defined for both
\hyperref[def:discrete-pmf]{discrete} and
\hyperref[def:continuous-pdf]{continuous} random variables.
\end{frameddefn}

\begin{framedprop}{Properties of the CDF (I)}\label{prop:cdf-properties-i}
A function $F(\cdot)$ mapping the real line to $[0,1]$ is a CDF
for some probability measure $P(\cdot)$ if and only if:
\begin{enumerate}
  \item $F(\cdot)$ is \emph{non-decreasing}: $x_1<x_2$ implies that
    $F(x_1)\leq F(x_2)$.
  \item $F(\cdot)$ is \emph{normalized}:
    \[
      \lim_{x\to-\infty}F(x)=0,
      \qquad
      \lim_{x\to+\infty}F(x)=1.
    \]
  \item $F(\cdot)$ is \emph{right-continuous}: for all $x$ we have
    \[
      F(x)=F(x^+),
      \qquad F(x^+):=\lim_{t\downarrow x}F(t).
    \]
\end{enumerate}
\end{framedprop}

\begin{framedprop}{Properties of the CDF (II/a)}
Let $F_X(\cdot)$ be the CDF of a random variable $X$. Then:
\begin{enumerate}
  \item For $x_1<x_2$,
    \[
      P(x_1<X\leq x_2)=P\bigl(X\in(x_1,x_2]\bigr)
        =F_X(x_2)-F_X(x_1).
    \]
    In the discrete case, be careful with the strict and non-strict
    inequalities.
  \item
    \[
      P(X>x)=P\bigl(X\in(x,+\infty)\bigr)=1-F_X(x).
    \]
  \item If $X$ is continuous with PDF $f_X(\cdot)$, then
    \[
      F_X(x)=\int_{-\infty}^{x}f_X(t)\,dt.
    \]
    A PDF can be recovered from the CDF as
    \[
      f_X(x)=F_X'(x)
    \]
    at points where $F_X$ is differentiable.
\end{enumerate}
\end{framedprop}

\begin{framedprop}{Properties of the CDF (II/b)}
Let $F_X(\cdot)$ be the CDF of a random variable $X$. Then:
\begin{enumerate}
  \item
    \[
      P(x_1<X\leq x_2)=P\bigl(X\in(x_1,x_2]\bigr)
        =F_X(x_2)-F_X(x_1),\qquad x_1<x_2.
    \]
  \item
    \[
      P(X>x)=P\bigl(X\in(x,+\infty)\bigr)=1-F_X(x).
    \]
  \item If $X$ is discrete with PMF $p_X(\cdot)$ and support
    $\mathcal X=\{x_1,x_2,\ldots\}$, then
    \[
      F_X(x)=\sum_{x_i\leq x}p_X(x_i)=P(X<x)+p_X(x).
    \]
    Hence
    \[
      p_X(x)=F_X(x)-F_X(x^-),
      \qquad F_X(x^-):=\lim_{t\uparrow x}F_X(t).
    \]
\end{enumerate}
\end{framedprop}

\begin{frameddefn}{Quantile function: continuous case}
\label{def:quantile-continuous}
Let $X$ be a continuous random variable with CDF $F_X$. For a
probability level $0<p<1$, its \emph{quantile function} is
\[
  Q(p):=\inf\{x\in\mathbb R:F_X(x)=p\}.
\]
It gives the smallest value $x$ such that $F_X(x)=P(X\leq x)=p$.
We also write $q_p$ instead of $Q(p)$.
\end{frameddefn}

\begin{frameddefn}{Quantile function: general case}
\label{def:quantile-general}
Let $X$ be any random variable, discrete or continuous, with CDF
$F_X$. For a probability level $0<p<1$, its \emph{quantile function} is
\[
  Q(p):=\inf\{x\in\mathbb R:F_X(x)\geq p\}.
\]
It gives the smallest threshold $x$ such that
$P(X\leq x)$ is at least $p$. We also write $q_p$ instead of $Q(p)$.
\end{frameddefn}

\begin{frameddefn}{Uniform distribution}
\label{def:uniform-distribution}
A continuous random variable $X$ has a \emph{uniform distribution}
with parameters $a<b$, where $a,b\in\mathbb R$, if its PDF and CDF
are
\[
  f_X(x)=\frac{1}{b-a}\mathbf 1_{[a,b]}(x),
\]
\[
  F_X(x)=
  \begin{cases}
    0, & x<a,\\
    \dfrac{x-a}{b-a}, & a\leq x\leq b,\\
    1, & x>b.
  \end{cases}
\]
We denote this distribution by $\operatorname{Unif}(a,b)$ and write
$X\sim\operatorname{Unif}(a,b)$. Here $\operatorname{Unif}$
abbreviates \emph{uniform}, and $\sim$ means \emph{has distribution}.
\end{frameddefn}

\begin{frameddefn}{Exponential distribution}
We write $T\sim\operatorname{Exp}(\lambda)$, with rate $\lambda>0$,
when its PDF and CDF are
\[
  f_T(t)=\lambda e^{-\lambda t}\mathbf1_{[0,+\infty)}(t),
  \qquad
  F_T(t)=
  \begin{cases}
    0,&t<0,\\
    1-e^{-\lambda t},&t\geq0.
  \end{cases}
\]
\end{frameddefn}

\begin{frameddefn}{Bernoulli distribution}
For $0\leq p\leq1$, we write $X\sim\operatorname{Ber}(p)$ when
\[
  p_X(x)=
  \begin{cases}
    1-p,&x=0,\\p,&x=1,\\0,&\text{otherwise}.
  \end{cases}
\]
The parameter $p$ is the success probability. The two masses add
to $1$.
\end{frameddefn}

\begin{frameddefn}{Binomial distribution}
For an integer $n\geq1$ and $0<p<1$, $X\sim\operatorname{Bin}(n,p)$
has PMF
\[
  p_X(k)=\binom nk p^k(1-p)^{n-k},\qquad k=0,\ldots,n,
\]
and zero elsewhere. For $p=0$, $X=0$ with probability $1$;
for $p=1$, $X=n$ with probability $1$.
\end{frameddefn}

\begin{frameddefn}{Beta distribution}
\label{def:beta-distribution}
For $\alpha,\beta>0$, the density of
$P\sim\operatorname{Beta}(\alpha,\beta)$ is
\[
  \pi(p)=\frac{p^{\alpha-1}(1-p)^{\beta-1}}{B(\alpha,\beta)},
  \qquad 0<p<1,
\]
and it is zero elsewhere. The denominator is the normalizing
constant
\[
  B(\alpha,\beta)=\int_0^1
       t^{\alpha-1}(1-t)^{\beta-1}\,dt.
\]
\end{frameddefn}

\begingroup
\renewcommand{\thefootnote}{\fnsymbol{footnote}}
\footnotetext[4]{In the notes, $\Gamma(a)=\int_0^{+\infty}t^{a-1}e^{-t}\,dt$ for $a>0$.}
\endgroup
\nopagebreak
\begin{frameddefn}{Gamma distribution}
Using the Gamma function\begingroup\renewcommand{\thefootnote}{\fnsymbol{footnote}}\footnotemark[4]\endgroup\ defined above, we write
$T\sim\operatorname{Gamma}(a,\lambda)$ when its PDF is
\[
  f_T(t)=\frac{\lambda^a}{\Gamma(a)}
           t^{a-1}e^{-\lambda t}\mathbf1_{(0,+\infty)}(t).
\]
\end{frameddefn}

\notesheading{Random vectors}

\begin{frameddefn}{Random vector}
On a common probability space $(\Omega,\mathcal F,P)$, a
\textbf{random vector} is a vector of random variables
\[
  \underline{X}:\Omega\longrightarrow\mathbb R^p,
  \qquad
  \underline{X}(\omega)=
  \begin{pmatrix}
    X_1(\omega)\\
    \vdots\\
    X_p(\omega)
  \end{pmatrix}.
\]
The integer $p$ is the dimension of the vector. Each component
$X_j:\Omega\to\mathbb R$ is a random variable.
\end{frameddefn}

\begin{frameddefn}{Independence of random variables}
\label{def:independent-random-variables}
Two random variables $X$ and $Y$ are independent when
\[
  P(X\in I,\ Y\in J)=P(X\in I)P(Y\in J)
  \quad\text{for every pair of intervals }I,J.
\]
\end{frameddefn}

\begingroup
\renewcommand{\thefootnote}{\fnsymbol{footnote}}
\footnotetext[5]{Here $X$ and $Y$ are discrete, with supports $S_X$ and $S_Y$.}
\endgroup
\nopagebreak
\begin{frameddefn}{Joint probability mass function}
The \textbf{joint PMF} assigns a probability to each pair of values\begingroup\renewcommand{\thefootnote}{\fnsymbol{footnote}}\footnotemark[5]\endgroup:
\[
  p_{X,Y}(x,y)=P(X=x,Y=y)
    =P\bigl(\{X=x\}\cap\{Y=y\}\bigr).
\]
It satisfies
\[
  p_{X,Y}(x,y)\geq0,
  \qquad
  \sum_{x\in S_X}\sum_{y\in S_Y}p_{X,Y}(x,y)=1.
\]
Some pairs in $S_X\times S_Y$ may have probability zero.
The joint PMF completely determines the distribution of $(X,Y)$.
\end{frameddefn}

\begin{framedprop}{Marginal PMFs}
The marginal distributions are obtained by summing out the
other variable:
\[
  p_X(x)=\sum_{y\in S_Y}p_{X,Y}(x,y),
  \qquad
  p_Y(y)=\sum_{x\in S_X}p_{X,Y}(x,y).
\]
\end{framedprop}

\begin{frameddefn}{Conditional PMFs}
For a value $y$ with $p_Y(y)>0$,
\[
  p_{X\mid Y}(x\mid y)
    =P(X=x\mid Y=y)
    =\frac{p_{X,Y}(x,y)}{p_Y(y)}.
\]
Similarly, for $x$ with $p_X(x)>0$,
\[
  p_{Y\mid X}(y\mid x)
    =P(Y=y\mid X=x)
    =\frac{p_{X,Y}(x,y)}{p_X(x)}.
\]
For each fixed conditioning value, the conditional probabilities
are nonnegative and sum to $1$ over the remaining variable.
\end{frameddefn}

\begin{frameddefn}{Independence of discrete random variables}
Two discrete random variables $X$ and $Y$ are independent if
and only if
\[
  p_{X,Y}(x,y)=p_X(x)p_Y(y)
  \qquad\text{for every }(x,y)\in S_X\times S_Y.
\]
Equivalently,
\[
  p_{X\mid Y}(x\mid y)=p_X(x)
  \quad\text{for every }x
  \text{ and every }y\text{ with }p_Y(y)>0.
\]
\end{frameddefn}

\begin{frameddefn}{Joint cumulative distribution function}
For any pair of random variables, discrete or otherwise,
\[
  F_{X,Y}(x,y)=P(X\leq x,Y\leq y),
  \qquad (x,y)\in\mathbb R^2.
\]
This is the probability of the lower-left region
$(-\infty,x]\times(-\infty,y]$ in the plane.
\end{frameddefn}

\begin{framedprop}{Marginal CDFs}
\label{prop:marginal-cdfs}
The marginal CDFs are obtained as limits of the joint CDF:
\[
  F_X(x)=\lim_{y\to+\infty}F_{X,Y}(x,y),
  \qquad
  F_Y(y)=\lim_{x\to+\infty}F_{X,Y}(x,y).
\]
\end{framedprop}

\begin{framedprop}{$X,Y$ independent $\Longleftrightarrow$ joint CDF $=F_XF_Y$}
\label{prop:independence-cdf}
Let $X$ and $Y$ be \textbf{arbitrary random variables}, with
CDFs $F_X$ and $F_Y$. They are \textbf{independent if and only if}
their joint CDF satisfies
\[
  F_{X,Y}(x,y)=F_X(x)F_Y(y)
  \qquad\text{for all }x,y\in\mathbb R.
\]
\end{framedprop}

\begin{frameddefn}{Joint PDF}
\label{def:joint-pdf}
A \textbf{joint probability density function}
$f_{\underline{X}}:\mathbb R^p\to[0,+\infty)$ satisfies
\[
  \int_{\mathbb R^p}f_{\underline{X}}(\underline{x})\,d\underline{x}=1,
\]
and gives the probability of a measurable region
$A\subseteq\mathbb R^p$ through
\[
  P(\underline{X}\in A)=\int_A f_{\underline{X}}(\underline{x})\,d\underline{x}.
\]
\end{frameddefn}

\begin{frameddefn}{Kernel of a PDF}
\label{def:pdf-kernel}
A \textbf{kernel of a PDF} is a nonnegative function
$g(\underline{x})$ proportional to that PDF. It specifies the
shape of the density, leaving a positive multiplicative constant
to be determined.

To obtain the PDF, compute the total integral
\[
  c=\int_{\mathbb R^p}g(\underline{x})\,d\underline{x}.
\]
If $0<c<+\infty$, then
\[
  f_{\underline{X}}(\underline{x})=\frac{g(\underline{x})}{c}
\]
is the corresponding PDF. The constant $c$ is the
\textbf{normalizing constant}.
\end{frameddefn}

\begin{framedprop}{Marginal PDFs}
\label{prop:marginal-pdfs}
Let $\underline{X}=(X_1,\ldots,X_p)^{\mathsf T}$ have joint PDF
$f_{\underline{X}}$. The density of any subset of components is
obtained by integrating out the remaining ones. Reordering the
components if necessary, the first $k<p$ components have density
\[
  f_{X_1,\ldots,X_k}(x_1,\ldots,x_k)
    =\int_{\mathbb R^{p-k}}
       f_{\underline{X}}(x_1,\ldots,x_p)\,dx_{k+1}\cdots dx_p.
\]
In particular, for $p=2$ with components $(X,Y)$,
\[
  f_X(x)=\int_{\mathbb R}f_{X,Y}(x,y)\,dy,
  \qquad
  f_Y(y)=\int_{\mathbb R}f_{X,Y}(x,y)\,dx.
\]
\end{framedprop}

\begin{frameddefn}{Conditional PDFs}
When a joint PDF exists, a conditional density of $X$ given $Y=y$
can be defined, for $f_Y(y)>0$, by
\[
  f_{X\mid Y}(x\mid y)=\frac{f_{X,Y}(x,y)}{f_Y(y)}.
\]
Similarly, for $f_X(x)>0$,
\[
  f_{Y\mid X}(y\mid x)=\frac{f_{X,Y}(x,y)}{f_X(x)}.
\]
The conditioning value is held fixed. Integrating the conditional
density over the other variable gives $1$.
\end{frameddefn}

\begin{framedprop}{$X,Y$ independent $\Longleftrightarrow$ joint PDF $=f_Xf_Y$}
\label{prop:independence-pdf}
Suppose $X$ and $Y$ \textbf{each have a PDF},
denoted by $f_X$ and $f_Y$. They are \textbf{independent if and only if}
their joint PDF exists and satisfies
\[
  f_{X,Y}(x,y)=f_X(x)f_Y(y)
  \qquad\text{almost everywhere}.
\]
\end{framedprop}

\begin{frameddefn}{Mixed joint function}
\label{def:mixed-joint-function}
Let $X\mid(P=p)$ have PMF $b(k\mid p)$ on $\{0,\ldots,n\}$, and let
$P$ have prior PDF $\pi(p)$ on $(0,1)$.
The \textbf{mixed joint function} of $(X,P)$ is
\[
  q(k,p)=b(k\mid p)\,\pi(p),
  \qquad k=0,\ldots,n,\quad 0<p<1.
\]
It gives the probability that $X$ equals $k$ and, at the same time,
$P$ falls in an interval $(a,b)$:
\[
  \Pr(X=k,\ a<P<b)=\int_a^b q(k,p)\,dp,
  \qquad 0\leq a<b\leq1.
\]
\end{frameddefn}

\begin{framedprop}{Marginal PMF in the mixed model}
\label{prop:mixed-marginal-pmf}
The marginal PMF of $X$, denoted by $m(k)$, is
\[
  m(k)=\Pr(X=k)=\int_0^1 b(k\mid p)\,\pi(p)\,dp,
  \qquad k=0,\ldots,n.
\]
\end{framedprop}

\begin{framedprop}{Posterior density in the mixed model}
Let $k$ be a count with $m(k)>0$. Given $X=k$, the latent variable
$P$ has the \textbf{posterior density}
\[
  f_{P\mid X}(p\mid k)=\frac{b(k\mid p)\,\pi(p)}{m(k)},
  \qquad 0<p<1.
\]
\end{framedprop}

\begin{frameddefn}{Bivariate copula}
A \textbf{copula} is a joint CDF of two random variables with
standard uniform marginals:
\[
  C(u,v)=P(U\leq u,V\leq v),
  \qquad U,V\sim\operatorname{Unif}(0,1).
\]
The variables $U$ and $V$ need not be independent. On $[0,1]^2$,
the marginal conditions are $C(u,1)=u$ and $C(1,v)=v$.
\end{frameddefn}

\begin{framedthm}{Sklar's theorem: bivariate form}
For any joint CDF $F_{X,Y}$ with marginals $F_X$ and $F_Y$,
there exists a copula $C$ such that
\[
  F_{X,Y}(x,y)=C\bigl(F_X(x),F_Y(y)\bigr).
\]
If both marginals are continuous, the copula is unique. Otherwise,
it is uniquely determined only on
$\operatorname{Ran}(F_X)\times\operatorname{Ran}(F_Y)$, the product
of the sets of values taken by the marginal CDFs.

Conversely, combining any two marginal CDFs with a copula using
this formula produces a joint CDF with those marginals.
\end{framedthm}

\begin{frameddefn}{Independent and identically distributed sample}
A random sample $(X_1,\ldots,X_n)$ is \textbf{IID} from $F_X$ if
\begin{itemize}
  \item the random variables $X_1,\ldots,X_n$ are mutually
    independent;
  \item each $X_i$ has the same CDF $F_X$.
\end{itemize}
Here $n$ is the sample size. Uppercase $X_i$ denotes a random
observation before sampling; lowercase $x_i$ denotes its observed
value. The realized dataset is $\mathcal D_n=(x_1,\ldots,x_n)$.
\end{frameddefn}

\begin{frameddefn}{Likelihood}
\label{def:likelihood}
For fixed observed data $\mathcal D_n=(x_1,\ldots,x_n)$ and
a model parameter $\theta\in\Theta$, the \textbf{IID likelihood} is
\[
  L(\theta;x_1,\ldots,x_n)=
  \begin{cases}
    \displaystyle\prod_{i=1}^n p_\theta(x_i),
       & \text{discrete model},\\[6pt]
    \displaystyle\prod_{i=1}^n f_\theta(x_i),
       & \text{model with a PDF}.
  \end{cases}
\]
\end{frameddefn}

\begin{frameddefn}{Maximum likelihood estimation}
\label{def:maximum-likelihood}
A \textbf{maximum likelihood estimator (MLE)} chooses
\[
  \widehat\theta_n\in
  \argmax_{\theta\in\Theta}L(\theta;x_1,\ldots,x_n),
\]
when a maximizer exists. Here $\widehat\theta_n$ is the parameter
estimate calculated from the $n$ observations, and $\argmax$
denotes the set of parameter values attaining the largest
likelihood.
\end{frameddefn}

\begin{frameddefn}{Empirical cumulative distribution function}
Given observations $x_1,\ldots,x_n$, the \textbf{empirical CDF}
is
\[
  \widehat F_n(t)=\frac1n\sum_{i=1}^n
    \mathbf 1_{\{x_i\leq t\}}.
\]
It is the fraction of observations at or below $t$.
The corresponding empirical distribution assigns mass $1/n$
to each observation, counting repetitions.
\end{frameddefn}

\notesheading{Stochastic simulations}

\begin{frameddefn}{Pseudorandom number generator}
\label{def:prng}
In a simplified description, a \emph{pseudorandom number generator}
(PRNG) uses a deterministic function
\[
  f:(0,1)\longrightarrow(0,1)
\]
and a starting value, or \emph{seed}, $u_0\in(0,1)$ to generate
\[
  u_0,\quad u_1=f(u_0),\quad u_2=f(f(u_0)),\quad
  \ldots,\quad u_n=f^n(u_0).
\]
The function is chosen so that the sequence imitates the statistical
behaviour of an IID $\operatorname{Uniform}(0,1)$ sample.
Here $f^n(u_0)$ means applying $f$ repeatedly $n$ times.
\end{frameddefn}

\begin{frameddefn}{Linear congruential generator}
\label{def:lcg}
A \emph{linear congruential generator} (LCG) uses integer parameters
$\alpha$, $\beta$, and $M>1$, and an integer seed
$x_0\in\{0,1,\ldots,M-1\}$. Its update rule is
\[
  x_{k+1}=f(x_k)=(\alpha+\beta x_k)\bmod M.
\]
Here $\bmod M$ gives the remainder after division by $M$, so every
state $x_k$ belongs to $\{0,1,\ldots,M-1\}$.
\end{frameddefn}

\begin{framedprop}{Quantile transform}
\label{prop:quantile-transform}
Let $F$ be any CDF and let
\[
  Q(p)=\inf\{x\in\mathbb R:F(x)\geq p\},\qquad 0<p<1,
\]
be its quantile function. If $U\sim\operatorname{Unif}(0,1)$ and
$Y=Q(U)$, then $Y$ has CDF $F$:
\[
  F_Y(y)=P(Y\leq y)=F(y),\qquad y\in\mathbb R.
\]
\end{framedprop}

\notesheading{Exploratory data analysis}

\begin{frameddefn}{Class of models}
In this density-estimation problem, a \textbf{class of models}
$\mathcal F$ is a set of candidate PDFs. Each $f\in\mathcal F$
represents one possible population distribution.
\end{frameddefn}

\begin{frameddefn}{Penalized log-likelihood}
\label{def:penalized-log-likelihood}
For fixed observed data $\mathcal D_n=(x_1,\ldots,x_n)$ from an
IID sample, consider a candidate PDF $f$. Let $J(f)\geq0$ be
its penalty, and let $\lambda>0$ be the weight given to that
penalty. The \textbf{penalized log-likelihood} is
\[
  \log L(f;\mathcal D_n)-\lambda J(f)
    =\sum_{i=1}^n\log f(x_i)-\lambda J(f).
\]
\end{frameddefn}

\clearpage
\section*{Relations}

\subsection*{From the PMF to the CDF}
\begin{tcolorbox}[colback=blue!3,colframe=blue!30!black,
  boxrule=0.4pt,arc=1mm,title={What we are trying to answer},
  colbacktitle=blue!10,coltitle=black,fonttitle=\bfseries]
If we know the probability of each individual
value of $X$, how do we calculate the probability that $X$ is at most $x$?
\end{tcolorbox}

For a discrete random variable, the event $X\leq x$ includes every
possible value $k$ that is at most $x$. These events $\{X=k\}$ are
pairwise disjoint, so we add their probabilities:
\[
  \boxed{\underbrace{F_X(x)}_{\text{CDF: to find}}=P(X\leq x)
    =\sum_{\substack{k\in\operatorname{supp}(X)\\k\leq x}}
      \underbrace{p_X(k)}_{\text{PMF: known}},
    \qquad x\in\mathbb R.}
\]
\subsection*{From the PDF to the CDF}
\label{rel:pdf-to-cdf}
\begin{tcolorbox}[colback=blue!3,colframe=blue!30!black,
  boxrule=0.4pt,arc=1mm,title={What we are trying to answer},
  colbacktitle=blue!10,coltitle=black,fonttitle=\bfseries]
We know how probability is spread over the possible values of $X$
(its PDF). How do we calculate the probability of getting a value
less than or equal to $x$?
\end{tcolorbox}
For a continuous random variable with PDF $f_X$, we calculate
$P(X\leq x)$ by integrating the PDF over $(-\infty,x]$:
\[
  \boxed{\underbrace{F_X(x)}_{\text{CDF: to find}}=P(X\leq x)
    =\int_{-\infty}^{x}\underbrace{f_X(t)}_{\text{PDF: known}}\,dt,
    \qquad x\in\mathbb R.}
\]
\clearpage
\subsection*{From the CDF to the PMF and PDF}
\begin{tcolorbox}[colback=blue!3,colframe=blue!30!black,
  boxrule=0.4pt,arc=1mm,title={What we are trying to answer},
  colbacktitle=blue!10,coltitle=black,fonttitle=\bfseries]
We know the CDF $F_X$. How do we recover the PMF $p_X$ if $X$ is
discrete, or the PDF $f_X$ if $X$ is continuous?
\end{tcolorbox}

\textbf{What we know.} The CDF is defined by
\[
  F_X(x):=P(X\leq x).
\]
For every threshold $x$, it gives the probability of obtaining a
value less than or equal to $x$. We want to recover the PMF or the
PDF, depending on the type of $X$.

\textbf{Discrete case: recover the PMF.}
The PMF is defined by
\[
  p_X(x):=P(X=x).
\]
The CDF includes all values up to $x$. To isolate the probability
of the single value $x$, we subtract the probability of values
strictly less than $x$:
\[
  P(X=x)=P(X\leq x)-P(X<x).
\]
Since the left limit satisfies
\[
  F_X(x^-):=\lim_{t\uparrow x}F_X(t)=P(X<x),
\]
we obtain
\[
  \boxed{\underbrace{p_X(x)}_{\text{PMF: to find}}
    =\underbrace{F_X(x)-F_X(x^-)}_{\text{CDF: known}}.}
\]
Thus the PMF at $x$ is the size of the jump of the CDF at $x$.

\textbf{Continuous case: recover the PDF.}
A PDF is a non-negative function $f_X$ such that
\[
  P(a\leq X\leq b)=\int_a^b f_X(t)\,dt,
  \qquad \int_{\mathbb R}f_X(t)\,dt=1.
\]
In particular,
\[
  F_X(x)=P(X\leq x)=\int_{-\infty}^{x}f_X(t)\,dt.
\]
The CDF is therefore the integral of the PDF. To recover a PDF,
we differentiate the CDF:
\[
  \boxed{\underbrace{f_X(x)}_{\text{PDF: to find}}
    =\frac{d}{dx}\underbrace{F_X(x)}_{\text{CDF: known}}.}
\]
This recovers a PDF at points where the derivative exists, by the
fundamental theorem of calculus.


\end{document}
